% Generated draft: requires mathematical and source review.
\documentclass[11pt]{article}
\usepackage[margin=25mm]{geometry}
\usepackage{fontspec}
\setmainfont{Latin Modern Roman}
\usepackage{amsmath,amssymb,mathtools,unicode-math}
\setmathfont{Latin Modern Math}
\usepackage{enumitem,needspace,xurl,hyperref}
\setlength{\parindent}{0pt}
\setlength{\parskip}{5pt}
\setlength{\emergencystretch}{3em}
\newcommand{\N}{\mathbb{N}}
\newcommand{\Z}{\mathbb{Z}}
\newcommand{\Q}{\mathbb{Q}}
\newcommand{\R}{\mathbb{R}}
\newcommand{\C}{\mathbb{C}}
\newcommand{\F}{\mathbb{F}}
\newcommand{\A}{\mathbb{A}}
\newcommand{\PP}{\mathbb{P}}
\newcommand{\E}{\mathbb{E}}
\newcommand{\eps}{\varepsilon}
\newcommand{\dd}{\,\mathrm{d}}
\begin{document}
\section*{UnsolvedMath ID: 30006300}
\textbf{Source problem code:} catalogue entry
\subsection{Definitions and mathematical statement}
Let \(Y\subset X_\Sigma\) be a subvariety of the toric variety denoted by \(X_\Sigma\) in the supplied record, and let \(\mathcal O_Y\) denote its structure sheaf, regarded as an object on \(X_\Sigma\). The hypotheses on \(X_\Sigma\), the ground field, and the intended class of subvarieties are not specified. Write \(\operatorname{Coh}_{\mathrm{dg}}(X_\Sigma)\) for the dg category of coherent sheaves intended by the source; its precise derived and enhancement conventions remain unspecified. For a complex algebra \(A\), write \(\operatorname{Perf}(A)\) for the category of perfect complexes of \(A\)-modules, with the categorical enhancement to be specified.

The supplied Question 6 asks whether, for a large class of such subvarieties \(Y\), one can construct the following data:
\begin{enumerate}
\item A deformation of categories \((\mathcal C_t)_t\), with distinguished special fiber \(\mathcal C_0\), whose general fiber is isomorphic to \(\operatorname{Coh}_{\mathrm{dg}}(X_\Sigma)\) in the intended categorical sense.
\item An object \(E_0\in\operatorname{Ob}(\mathcal C_0)\) that deforms to \(\mathcal O_Y\) under the identification of the general fiber. Here \(E_0\) denotes the object written as \(C_0\) in the source, to distinguish it from the special-fiber category.
\item A possibly singular nonpositively curved space \(L_V\).
\item A noncommutative neighborhood of \(E_0\) described by a functor of the displayed form
\[
F\colon\mathcal C_0\longrightarrow
\operatorname{Perf}\!\left(\mathbb C[\pi_1(L_V)]\right),
\]
where \(\mathbb C[\pi_1(L_V)]\) denotes the complex group algebra of the fundamental group of \(L_V\). The source permits an appropriate modification of the target category to account for singularities of \(L_V\).
\end{enumerate}

This is a programmatic existence question with unresolved specifications. The deformation parameter space, the meaning of deformation of categories and objects, the categorical meaning of isomorphism, the curvature condition for singular spaces, and the properties making the displayed functor a noncommutative neighborhood are not supplied. The original wording uses a general fiber, whereas the cleaned statement uses a generic fiber; their relationship requires a specified parameter space. The symbol \(V\), the fundamental-group basepoint and connectivity conventions, and the permitted singularity modifications are also unspecified. The displayed functor follows the category-domain interpretation of the source notation; the source does not clarify whether its domain should instead be a subcategory associated with \(E_0\).

The source offers the Donaldson divisor theorem in symplectic geometry as motivation for the suggestion that noncommutative neighborhoods may arise through deformations. It supplies neither a precise version of that theorem nor a deduction establishing the requested construction.
\subsection{Short English statement}
For a large, unspecified class of subvarieties \(Y\subset X_\Sigma\), the question asks for a categorical deformation recovering the coherent-sheaf dg category, a special-fiber object deforming to \(\mathcal O_Y\), a possibly singular nonpositively curved space, and a functor describing a noncommutative neighborhood through perfect complexes over its fundamental-group algebra. Several essential definitions and hypotheses remain unspecified, so this is a statement-only draft of a programmatic question requiring human mathematical review.
\subsection{Sources}
\begin{itemize}
\item Supplied local UnsolvedMath registry record, permanent ID \(30006300\), title \emph{Categorical Noncommutative Neighborhoods of Toric Subvarieties}, exact source code \texttt{catalogue entry}, Question 6 in the Oberwolfach Reports Open Problems collection. No individual proposer is identified in the supplied record; the proposed year is 2025.
\item Supplied report citation and literature-source link: \emph{Toric Geometry} (2025), \emph{Oberwolfach Reports}, pp.\ 883--940, DOI \texttt{10.4171/owr/2025/19}. \url{https://doi.org/10.4171/owr/2025/19}. The record describes Oberwolfach Reports as documenting mathematical workshops and their open-problem discussions. The report contents have not been independently consulted for this draft.
\item Supplied literature assessment, checked on \texttt{2026-08-22}: the record labels the problem open and describes the construction as programmatic, reporting that no general theorem resolving it was located in the supplied triage. This is supplied dated curation, not an independently verified literature assessment. The listed source checked is the Oberwolfach report linked above; time-sensitive status and cited claims require independent verification.
\item The supplied question invokes the Donaldson divisor theorem in symplectic geometry as motivation. No separate bibliographic reference or precise theorem statement is supplied for that invocation.
\end{itemize}
\end{document}
